\section*{Capítulo \ref{ch:b3:stem-cells} --- Células madre, regeneración y envejecimiento}

\begin{solution}{exo:b3:stem-cells:1}
\omterm{def:b3:stem-cells:stem}{Autorrenovación}: una división que produce al menos una hija que sigue
siendo \omterm{def:b3:stem-cells:stem}{célula madre} ---las células \emph{Lgr5} de la base de la cripta.
Potencia: el abanico de destinos que una célula todavía puede adoptar
---la \omterm{def:b3:stem-cells:stem}{célula madre} de la cripta es multipotente y hace células
absortivas, caliciformes, enteroendocrinas, de Paneth y en penacho.
\omterm{def:b3:stem-cells:stem}{Nicho}: las \omterm{prop:b3:stem-cells:crypt}{células de Paneth} y sus señales Wnt, Notch y EGF, que
mantienen en su sitio el carácter de madre. \omterm{def:b3:stem-cells:stem}{Célula amplificadora de
tránsito}: una hija comprometida que se divide cuatro o cinco veces más
en la cripta antes de diferenciarse camino de la vellosidad.
\end{solution}

\begin{solution}{exo:b3:stem-cells:2}
Células aisladas de médula produjeron colonias esplénicas que contenían
glóbulos rojos, granulocitos y plaquetas: una célula, varios linajes
---multipotencia---, y la respuesta lineal a la dosis mostró que cada
colonia procedía de una sola célula. La \omterm{def:b3:stem-cells:stem}{autorrenovación} necesitó un
segundo ratón: células de una colonia, inyectadas en otro receptor
irradiado, hicieron colonias nuevas, de modo que la célula fundadora
había producido hijas con su misma capacidad.
\end{solution}

\begin{solution}{exo:b3:stem-cells:3}
\emph{Oct4}, \emph{Sox2}, \emph{Klf4}, \emph{c-Myc}. La \omterm{def:b3:chromatin-epigenetics:nucleosome}{cromatina} del
fibroblasto ha cerrado los genes de la \omterm{def:b3:stem-cells:pluripotency}{pluripotencia} tras
\omterm{def:b3:chromatin-epigenetics:nucleosome}{heterocromatina} y \omterm{def:b3:chromatin-epigenetics:methylation}{metilación del ADN}, de modo que los factores
encuentran al principio pocos sitios accesibles y han de reclutar
enzimas remodeladoras a lo largo de muchos \omterm{def:b3:cell-cycle-apoptosis:cdk}{ciclos celulares}; la célula
ha de apagar además su propio programa, escapar de la \omterm{def:b3:cell-cycle-apoptosis:other}{senescencia} y
atravesar un estado intermedio estocástico en el que la mayoría se
atasca o muere. Solo la rara célula en la que todos los pasos salen
bien emerge pluripotente, semanas después.
\end{solution}

\begin{solution}{exo:b3:stem-cells:4}
La planaria conserva \omterm{def:b3:stem-cells:stem}{células madre} adultas pluripotentes, los
\omterm{def:b3:stem-cells:regeneration}{neoblastos}, que migran a la herida y construyen lo que falte. La
salamandra no tiene una reserva así: las células diferenciadas del
muñón se desdiferencian en \omterm{def:b3:stem-cells:stem}{progenitores} que, con las \omterm{def:b3:stem-cells:stem}{células madre}
residentes del tejido, forman un \omterm{def:b3:stem-cells:regeneration}{blastema} que vuelve a ejecutar el
programa de la extremidad. Los injertos de tejidos marcados de uno en
uno de Kragl mostraron que las células del \omterm{def:b3:stem-cells:regeneration}{blastema} se mantienen fieles
a su origen ---el músculo hace músculo y el cartílago, cartílago---, de
modo que el \omterm{def:b3:stem-cells:regeneration}{blastema} es una mezcla de \omterm{def:b3:stem-cells:stem}{progenitores} restringidos a su
linaje y no una masa pluripotente.
\end{solution}

\begin{solution}{exo:b3:stem-cells:5}
$(11 - 5)/0.075 = 80$ duplicaciones. Desde \qty{8}{kb}:
$(8 - 5)/0.075 = 40$. Con una compensación del \qty{60}{\%} la pérdida
neta es de \qty{30}{bp}: $6000/30
= 200$.
\end{solution}

\begin{solution}{exo:b3:stem-cells:6}
Catorce divisiones de \omterm{def:b3:stem-cells:stem}{células madre} al día rinden 28 hijas; 14 han de
seguir siendo \omterm{def:b3:stem-cells:stem}{células madre}, de modo que 14 se comprometen ---una hija
comprometida por \omterm{def:b3:stem-cells:stem}{célula madre} y día. Cada una se compromete a
$2^{4} = 16$ células: $14\times 16 = 224$ células de vellosidad al día,
el orden de las \num{300} observadas (una quinta división de tránsito
da \num{448}). El piso de las \omterm{def:b3:stem-cells:stem}{células madre} aporta divisiones por valor
de catorce células, y el de tránsito, las otras \num{200}.
\end{solution}

\begin{solution}{exo:b3:stem-cells:7}
$2\times 10^{11}/2^{20} = 1.9\times 10^{5}$ hijas comprometidas al día;
con \num{10000} \omterm{def:b3:stem-cells:stem}{células madre}, 19 por \omterm{def:b3:stem-cells:stem}{célula madre} y día ---frente a
una división al año. El recuento de veinte divisiones es por linaje,
desde la \omterm{def:b3:stem-cells:stem}{célula madre} hasta el glóbulo rojo, pero los \omterm{def:b3:stem-cells:stem}{progenitores}
intermedios no son de un solo uso: las reservas de \omterm{def:b3:stem-cells:stem}{progenitores}
multipotentes y comprometidos se sostienen durante semanas con sus
propias divisiones, de modo que casi toda la división ocurre en los
pisos de \omterm{def:b3:stem-cells:stem}{progenitores} y a la \omterm{def:b3:stem-cells:stem}{célula madre} se la llama rara vez.
\end{solution}

\begin{solution}{exo:b3:stem-cells:8}
$\mu(50) = 10^{-3}\mathrm{e}^{1.74} = 5.7\times 10^{-3}$; $\mu(70) =
10^{-3}\mathrm{e}^{3.48} = 0.032$; $\mu(90) = 10^{-3}\mathrm{e}^{5.22} =
0.19$ al año. Con $\mu_{0}/\gamma = 0.0115$: $S(70) = \exp[-0.0115
\times 31.5] = 0.70$ y $S(90) = \exp[-0.0115\times 184] = 0.12$.
Mediana: $30 + \ln(61)/0.087 = 77$. Con $\mu_{0}$ a la mitad:
$\ln(121)/0.087 = 55$, mediana 85 ---se gana un tiempo de duplicación,
ocho años. Con $\gamma$ a la mitad, $0.0435$: $\ln(31)/0.0435 = 79$,
mediana 109: frenar el exponente gana cuatro veces más que reducir a la
mitad la constante.
\end{solution}

\begin{solution}{exo:b3:stem-cells:9}
El tercio restante ha de triplicarse: $\log_{2}3 = 1.6$ divisiones por
célula. \omterm{def:b3:stem-cells:regeneration}{Crecimiento compensador}, porque los lóbulos existentes se
agrandan por división de células en su sitio; los lóbulos retirados,
con sus conductos, sus vasos y su arquitectura, no se reconstruyen ---el
hígado recupera su masa y su función, pero no su forma.
\end{solution}

\begin{solution}{exo:b3:stem-cells:10}
En cada suceso el tamaño $k$ del clon pasa a $k + 1$ o a $k - 1$ con
igual probabilidad, de modo que su tamaño esperado no cambia nunca;
cuando el paseo termina vale $N$ (se impone) o $0$ (se pierde), y la
esperanza $N\,P + 0 = k$ da $P = k/N$. Una sola célula: $1/14$. Una
mutación neutra en una \omterm{def:b3:stem-cells:stem}{célula madre} se fija en su cripta con
probabilidad $1/14$ y en los demás casos se pierde en unos meses ---la
deriva purga trece de cada catorce mutaciones.
\end{solution}

\begin{solution}{exo:b3:stem-cells:11}
El paseo está ahora sesgado: el clon gana una plaza más a menudo de lo
que pierde otra, de modo que su tamaño esperado crece y la probabilidad
de imponerse se acerca a la certeza a medida que $k$ crece ---para una
ventaja del \qty{10}{\%} en un \omterm{def:b3:stem-cells:stem}{nicho} de catorce, alrededor del triple
de $1/14$ para una sola célula, y todavía más una vez que el clon ocupa
unas cuantas plazas. A lo largo de décadas, las \omterm{def:b3:stem-cells:stem}{células madre} con esas
mutaciones (\emph{DNMT3A}, \emph{TET2} en la médula) se apoderan de sus
\omterm{def:b3:stem-cells:stem}{nichos}, y a los setenta años una de cada cinco personas tiene una parte
grande de su sangre procedente de un solo clon. No es \omterm{def:b3:cancer-biology:hallmarks}{cáncer}: las
células siguen diferenciándose y obedeciendo a su jerarquía. Pero la
mutación siguiente tiene ya un clon de miles de millones en el que
surgir, y el riesgo de leucemia se multiplica por diez.
\end{solution}

\begin{solution}{exo:b3:stem-cells:12}
Mediana $= 30 + \gamma^{-1}\ln(1 + \gamma\ln 2/\mu_{0})$: para $\mu_{0} =
10^{-2}$, $30 + \ln 7/0.087 = 52$; para $10^{-3}$, 77; para $10^{-4}$, $30 +
\ln 604/0.087 = 104$. En la forma pura de Gompertz, cada reducción de
diez veces añade los mismos $\ln 10/\gamma = 26$ años. Los rendimientos
decrecen en la práctica porque las ganancias históricas vinieron de
retirar un término independiente de la edad ---infección, parto,
accidente, la constante $A$ de Makeham en $\mu = A
+ \mu_{0}\mathrm{e}^{\gamma t}$--- y, una vez que $A$ es pequeño frente
a la exponencial en las edades que importan, recortarlo más no cambia
casi nada; el $\mu_{0}$ intrínseco se ha movido mucho menos. Una
reducción del \qty{20}{\%} de $\gamma$ (hasta $0.070$) da
$30 + \ln 49/0.070 = 86$: nueve años, ganados a todas las edades y no
por aplazar unas pocas causas de muerte.
\end{solution}

\begin{solution}{pb:b3:stem-cells:1}
\textbf{1.} $3.5\times 10^{11}$ al día, $4\times 10^{6}$ por segundo.
\textbf{2.} $2^{20} \approx 10^{6}$ células por hija comprometida;
$3.5\times 10^{11}/10^{6} = 3.3\times 10^{5}$ hijas comprometidas al día.
\textbf{3.} $33$ por \omterm{def:b3:stem-cells:stem}{célula madre} y día, \num{12000} al año.
\textbf{4.} Las reservas de \omterm{def:b3:stem-cells:stem}{progenitores} han de sostenerse solas: los
\omterm{def:b3:stem-cells:stem}{progenitores} multipotentes y comprometidos se dividen durante semanas y
mantienen sus propios números, de modo que casi todas las veinte
divisiones ocurren en pisos que se renuevan a sí mismos, y la \omterm{def:b3:stem-cells:stem}{célula
madre} aporta una hija comprometida rara vez ---del orden de una vez al
año.
\textbf{5.} $10 - 20\times 0.06 = \qty{8.8}{kb}$: muy por encima de
$L_{\text{crit}}$, y el glóbulo rojo no vuelve a dividirse nunca; la
reserva solo importa para los \omterm{def:b3:stem-cells:stem}{progenitores} que se reutilizan.
\textbf{6.} Pérdida neta $0.3\times 60 = \qty{18}{bp}$; $6000/18 = 333$
divisiones; a una al año, $333$ años ---mucho más allá de una vida, de
modo que los telómeros no limitan a la propia \omterm{def:b3:stem-cells:stem}{célula madre}; otros daños
sí.
\textbf{7.} $14$ divisiones al día; $7$ reemplazan \omterm{def:b3:stem-cells:stem}{células madre}
perdidas y $7$ se comprometen; $7\times 2^{4} = 112$ células de
vellosidad al día.
\textbf{8.} $N^{2} = 196$ sucesos por célula a uno al día: unos
\qty{200}{d}, de seis a siete meses ---el tiempo en el que las criptas
del confeti pasaron a tener un solo color.
\textbf{9.} $1/14 \approx \qty{7}{\%}$; $10^{7}/14 \approx 7\times 10^{5}$
criptas.
\textbf{10.} La deriva descarta trece de cada catorce mutaciones en unos
meses. Una \omterm{def:b3:stem-cells:stem}{división asimétrica} de por vida conservaría toda \omterm{def:b3:stem-cells:stem}{célula madre}
mutante durante toda la vida, de modo que las mutaciones se acumularían
en cada linaje sin purgarse jamás.
\textbf{11.} El carácter de madre es una posición en el \omterm{def:b3:stem-cells:stem}{nicho} y un
conjunto de señales, no una propiedad permanente de una célula:
cualquier célula que alcance las señales de Paneth puede asumir el
papel.
\textbf{12.} Indúzcase una marca multicolor en las \omterm{def:b3:stem-cells:stem}{células madre} en un
instante y síganse las criptas. Deriva simétrica: las criptas pasan a
tener un solo color en unos meses. \omterm{def:b3:stem-cells:stem}{División asimétrica}: cada cripta
conserva todos sus colores indefinidamente, cada uno como una cinta
fina que sube por la vellosidad.
\textbf{13.} $6000/60 = 100$ duplicaciones.
\textbf{14.} Le quedan unas $75$: la célula cuenta divisiones y no
tiempo de calendario, ya que la década en el congelador no le costó
nada.
\textbf{15.} $(8000 - 4000)/30 = 133$ años después de los veinte ---a
los 153. La longitud media del telómero de los leucocitos no limita por
sí misma la vida; importan más los telómeros más cortos y las demás
marcas.
\textbf{16.} La \omterm{def:b3:cancer-biology:telomerase}{telomerasa} retira una barrera, la \omterm{thm:b3:stem-cells:hayflick}{senescencia
replicativa}, y por eso se encuentra en casi todos los \omterm{def:b3:cancer-biology:hallmarks}{cánceres}, que
necesitan dividirse sin límite; pero los fibroblastos a los que se dio
\omterm{def:b3:cancer-biology:telomerase}{telomerasa} conservaron sus puntos de control, la inhibición por
contacto y una p53 intacta, de modo que la división ilimitada por sí
sola no los volvió \omterm{def:b3:cancer-biology:hallmarks}{cánceres}. Necesaria, no suficiente.
\textbf{17.} $\mu(40) = 2.4\times 10^{-3}$, $\mu(60) = 0.014$, $\mu(80) =
0.077$ y $\mu(100) = 0.44$ al año. $S = \exp[-0.0115(\mathrm{e}^{\gamma
(t-30)} - 1)]$: $0.98$ a los 40, $0.87$ a los 60, $0.42$ a los 80 y
$0.006$ a los 100.
\textbf{18.} Mediana $30 + \ln 61/0.087 = 77$. Sobrevive el diez por
ciento cuando $\mathrm{e}^{\gamma(t-30)} = 1 + \gamma\ln 10/\mu_{0} = 201$: $t = 30 +
5.3/0.087 = 91$.
\textbf{19.} \qty{0.67}{mm/d}; $\log_{2}1000 = 10$ duplicaciones en 60
días, una cada seis días.
\textbf{20.} Unas $10^{8}$ divisiones, de modo que $0.1$ mutaciones
oncogénicas por \omterm{def:b3:stem-cells:regeneration}{regeneración} ---una de cada diez extremidades. Los
ajolotes casi nunca desarrollan \omterm{def:b3:cancer-biology:hallmarks}{tumores}: son de sangre fría y lentos,
con una supresión tumoral robusta, y las células del \omterm{def:b3:stem-cells:regeneration}{blastema} siguen
restringidas a su linaje y se rediferencian en lugar de seguir
dividiéndose; el tejido en \omterm{def:b3:stem-cells:regeneration}{regeneración} incluso suprime los \omterm{def:b3:cancer-biology:hallmarks}{tumores}
inducidos. Un mamífero tiene más células, una temperatura y un ritmo
metabólico mayores y una vida más larga a lo largo de la cual podría
crecer un clon escapado.
\textbf{21.} Desnérvese el muñón e implántese una bolita que libere la
proteína candidata, o injértese tejido acondicionado por nervio: si la
\omterm{def:b3:stem-cells:regeneration}{regeneración} se restablece, hay un factor difusible. Candidatos
identificados en tritones y ajolotes: la proteína de gradiente anterior
nAG, la neurregulina-1 y los FGF.
\textbf{22.} La \omterm{prop:b3:stem-cells:crypt}{señalización Wnt} es alta en la cola y baja en la cabeza,
y la herida del extremo anterior de un trozo hace normalmente una
cabeza porque allí la Wnt es baja; silenciar un inhibidor de Wnt eleva
la Wnt por todas partes, ambas heridas leen ``posterior'' y crecen dos
colas. El trozo conoce su propio eje por el gradiente residual de su
musculatura, y la herida vuelve a establecer el gradiente respecto de
él.
\textbf{23.} El tercio restante ha de triplicarse: todo hepatocito se
divide una vez (hasta dos tercios) y la mitad se divide otra vez ---una
media de $1.6$ divisiones. Los lóbulos supervivientes se hinchan hasta
la masa original; los lóbulos que faltan, con sus conductos y sus
vasos, no se reconstruyen.
\textbf{24.} Una \omterm{def:b3:stem-cells:regeneration}{regeneración} que sustituyera el tejido gastado y
dañado frenaría la acumulación que mide el exponente, bajaría $\gamma$
y aplanaría la curva ---los ajolotes muestran poca \omterm{def:b3:cell-cycle-apoptosis:other}{senescencia}. El
precio sería un término constante mayor por los \omterm{def:b3:cancer-biology:hallmarks}{cánceres} en un cuerpo
cálido, grande y longevo.
\textbf{25.} Recuento de Hayflick de unas $100$ duplicaciones; una
\omterm{def:b3:stem-cells:stem}{célula madre} asistida por \omterm{def:b3:cancer-biology:telomerase}{telomerasa} tiene unas $330$ divisiones, más
de tres siglos a una al año; vida mediana de $77$ años.
\end{solution}
